% einstein.tex — Paper 6, Section V: Einstein's Equations
% ASSERT_CONVENTION: natural_units=natural, metric_signature=mostly_plus, entropy_base=nats, modular_hamiltonian=K_minus_ln_rho, coupling_convention=H_sum_hxy, state_normalization=Tr_rho_1
\section{Einstein's Equations}
\label{sec:einstein}

We present two independent routes from the self-modeling lattice to
Einstein's field equations, with complementary strengths.  Route~A
(Sec.~\ref{sec:route-a}) applies Jacobson's (2016) entanglement
equilibrium argument~\cite{Jacobson2016} using the conformal modular
Hamiltonian; it is exact in the conformal regime ($d = 1$ WZW, or the
window $a \ll \ell \ll 1/m$ in higher $d$).  Route~B
(Sec.~\ref{sec:route-b}) uses Lovelock's uniqueness
theorem~\cite{Lovelock1971, Lovelock1972} to show that \emph{any}
tensorial equation relating entanglement perturbations to geometry
perturbations must be Einstein's equation in $d \geq 2$ spatial
dimensions.  Route~B does not require a conformal modular Hamiltonian.

\subsection{Route A: Jacobson entanglement equilibrium (conformal regime)}
\label{sec:route-a}

\subsubsection{Entropy decomposition}
\label{sec:setup}

Consider a point $p$ in the emergent spacetime manifold $(M, \gab)$
obtained from the Wilsonian continuum limit
(Sec.~\ref{sec:continuum}).  Choose Riemann normal coordinates at $p$
so that $\gab(p) = \eta_{ab}$ and the Christoffel symbols vanish.  Let
$B$ be a geodesic ball of proper radius $\ell$ centered at $p$ in the
spacelike hypersurface $\Sigma = \{x^0 = 0\}$, with $a \ll \ell \ll
L_{\mathrm{curv}}$ (the curvature radius).  The causal diamond
$\mathcal{D}(B)$ is the domain of dependence of~$B$.

The entanglement entropy of ball $B$ decomposes as~\cite{Jacobson2016}
\begin{equation}
  \ent = \ent_{\mathrm{UV}} + \ent_{\mathrm{mat}},
  \label{eq:s-decomp}
\end{equation}
where $\ent_{\mathrm{UV}} = \eta\, \mathcal{A}$ is the
area-proportional UV contribution (with $\eta = 1/(4\Gnewton)$ the
entropy density per unit area and $\mathcal{A}$ the area of
$\partial B$), and $\ent_{\mathrm{mat}}$ is the finite,
state-dependent matter contribution.  The UV entropy density $\eta$
is a constant of the lattice structure; it does not depend on the
state.

Under a first-order perturbation away from the maximally symmetric
spacetime (MSS) vacuum:
\begin{equation}
  \delta\ent = \delta\ent_{\mathrm{UV}} + \delta\ent_{\mathrm{mat}}
  = \eta\, \delta\mathcal{A} + \delta\ent_{\mathrm{mat}}.
  \label{eq:ds-decomp}
\end{equation}

\subsubsection{Geometric variation (Raychaudhuri)}
\label{sec:raychaudhuri}

The change in area $\delta\mathcal{A}$ of the entangling surface
$\partial B$ due to the curvature perturbation is computed from the
Raychaudhuri equation for the null generators of the causal
diamond~\cite{Wald1984, Jacobson2016}.  For a twist-free null geodesic
congruence with tangent $k^a$ ($k^a k_a = 0$), the expansion
$\theta = \nabla_a k^a$ satisfies
\begin{equation}
  \frac{d\theta}{d\lambda}
  = -\frac{1}{d-1}\,\theta^2 - \sigma_{ab}\sigma^{ab} - \Rab\, k^a k^b,
  \label{eq:raychaudhuri}
\end{equation}
where $d$ is the number of spatial dimensions and $\lambda$ is the
affine parameter.

At the bifurcation surface $\partial B$, the unperturbed expansion and
shear vanish ($\theta = \sigma_{ab} = 0$).  To first order in the
curvature perturbation, $\theta^2$ and $\sigma^2$ are second order, so
\begin{equation}
  \frac{d\theta^{(1)}}{d\lambda} = -\Rab^{(1)}\, k^a k^b.
  \label{eq:ray-linear}
\end{equation}
Integrating over the null congruence emanating from the future tip of
$\mathcal{D}(B)$ and performing the angular average over null
directions~\cite{Jacobson2016}, the area variation in the small-ball
limit ($\ell \ll L_{\mathrm{curv}}$) is
\begin{equation}
  \delta\mathcal{A}
  = -\frac{\Omega_{d-1}\, \ell^{d+1}}{2d}
  \left[(d+1)\, \Rab\, n^a n^b + \mathcal{R}\right]
  + O(\ell^{d+3}),
  \label{eq:delta-area}
\end{equation}
where $\Omega_{d-1} = 2\pi^{d/2}/\Gamma(d/2)$ is the area of the unit
$(d{-}1)$-sphere, $n^a$ is the unit timelike normal to $\Sigma$, and
$\mathcal{R} = g^{ab}\Rab$ is the Ricci scalar.

The UV entropy variation is therefore
\begin{equation}
  \delta\ent_{\mathrm{UV}} = \eta\,\delta\mathcal{A}
  = -\frac{\eta\,\Omega_{d-1}\, \ell^{d+1}}{2d}
  \left[(d+1)\,\Rab\, n^a n^b + \mathcal{R}\right].
  \label{eq:ds-uv}
\end{equation}
When $\Rab\, n^a n^b > 0$ (i.e., the null energy condition holds),
$\delta\ent_{\mathrm{UV}} < 0$: positive energy causes focusing, which
decreases the area and hence the UV entropy.

For conformal matter, the Weyl tensor contribution to
$\delta\ent_{\mathrm{UV}}$ is exactly cancelled by a corresponding
contribution to $\delta\ent_{\mathrm{mat}}$~\cite{Jacobson2016}, so
only the Ricci part survives.

\subsubsection{Matter variation (conformal modular Hamiltonian)}
\label{sec:conformal-approx}

From the entanglement first law~\eqref{eq:first-law},
$\delta\ent_{\mathrm{mat}} = \delta\langle{\modH}_B\rangle$.  For the
vacuum state of a conformal field theory restricted to a geodesic ball
$B$ of radius $\ell$, the Casini--Huerta--Myers (CHM) modular Hamiltonian
takes the exact form~\cite{CasiniHuertaMyers2011}
\begin{equation}
  {\modH}_B = 2\pi \int_B d^d\!x\; \frac{\ell^2 - |\mathbf{x}|^2}{2\ell}\; T_{00}(\mathbf{x}).
  \label{eq:chm}
\end{equation}
Under a perturbation that produces a stress-energy tensor~$\Tab$, in
the small-ball limit where $T_{00}$ is approximately constant over $B$:
\begin{equation}
  \delta\ent_{\mathrm{mat}}
  = \frac{2\pi\, \Omega_{d-1}}{d(d+2)}\, \ell^{d+1}\, \Tab\, n^a n^b.
  \label{eq:ds-mat}
\end{equation}
For positive energy density ($\Tab\, n^a n^b > 0$),
$\delta\ent_{\mathrm{mat}} > 0$: matter excitations increase the
matter sector entropy.

\emph{Conformal approximation.}---The CHM formula~\eqref{eq:chm} is
exact for conformal field theories in causal
diamonds~\cite{CasiniHuertaMyers2011}.  For nonconformal fields,
corrections of order $O\!\left((m\ell)^{2\Delta}\right)$ appear, where
$m$ is the mass scale breaking conformal invariance and $\Delta$ is
its scaling dimension~\cite{Speranza2016}.  In the window
$a \ll \ell \ll 1/m$ (if it exists), these corrections are suppressed.
For the one-dimensional AFM Heisenberg chain ($n = 2$, $J > 0$),
which flows to the $SU(2)_1$ WZW CFT with $c = 1$, the theory is
exactly conformal and this approximation introduces no error.
Jacobson~\cite{Jacobson2016} conjectures that the leading-order
Einstein equation survives nonconformal corrections; this is
physically motivated but not proven.

\subsubsection{Entanglement equilibrium implies Einstein's equation}
\label{sec:ee-implies-einstein}

The MVEH---identified as the only available vacuum selection criterion
in a pre-geometric framework
(Sec.~\ref{sec:mveh})---requires entanglement equilibrium:
\begin{equation}
  \delta\ent = \delta\ent_{\mathrm{UV}} + \delta\ent_{\mathrm{mat}} = 0.
  \label{eq:ee}
\end{equation}
Substituting Eqs.~\eqref{eq:ds-uv} and~\eqref{eq:ds-mat} and
cancelling the common factor $\Omega_{d-1}\, \ell^{d+1}/d > 0$:
\begin{equation}
  \frac{\eta(d+1)}{2}\, \Rab\, n^a n^b + \frac{\eta\, \mathcal{R}}{2}
  = \frac{2\pi}{d+2}\, \Tab\, n^a n^b.
  \label{eq:scalar-ee}
\end{equation}
This must hold for \emph{all} unit timelike vectors $n^a$ at every
point.  A symmetric tensor equation
$A_{ab}\, n^a n^b = B_{ab}\, n^a n^b$ for all unit timelike $n^a$
implies $A_{ab} = B_{ab} + f\, \gab$ for some scalar~$f$ (the trace
freedom arises because $\gab\, n^a n^b = -1$ for all such $n^a$).

Extracting the tensor equation up to this trace freedom, converting to
the Einstein tensor $\Gab = \Rab - \frac{1}{2} \mathcal{R}\, \gab$, and
absorbing all $\gab$ terms (including the background MSS curvature)
into an undetermined cosmological constant~$\Lambda$
yields~\cite{Jacobson2016}
\begin{equation}
  \boxed{\Gab + \Lambda\, \gab = 8\pi\Gnewton\, \Tab}
  \label{eq:einstein}
\end{equation}
with
\begin{equation}
  \boxed{\Gnewton = \frac{1}{4\eta}.}
  \label{eq:newton}
\end{equation}
This is the standard Einstein field equation.  The cosmological
constant $\Lambda$ is \emph{undetermined}: it is an integration
constant of the derivation, not a predicted
quantity~\cite{Jacobson2016, JacobsonSperanza2019}.

\emph{Sign chain.}---If the null energy condition (NEC) holds
($\Tab\, k^a k^b \geq 0$ for all null $k^a$), the derivation produces
attractive gravity via the following sign chain:
\begin{enumerate}
\item Positive energy density (NEC): $\Tab\, n^a n^b > 0$.
\item Matter entropy increases: $\delta\ent_{\mathrm{mat}} > 0$
  (Eq.~\ref{eq:ds-mat}).
\item Entanglement equilibrium: $\delta\ent = 0$ requires
  $\delta\ent_{\mathrm{UV}} < 0$.
\item UV entropy decreases: $\delta\ent_{\mathrm{UV}} < 0$ implies
  $\delta\mathcal{A} < 0$ (area decreases, since $\eta > 0$).
\item Raychaudhuri focusing: $\delta\mathcal{A} < 0$ implies
  $\Rab\, k^a k^b > 0$ (null geodesics converge).
\item Einstein equation: $\Gnewton > 0$ (since $\eta > 0$), and
  positive $\Tab$ produces positive $\Gab$, giving attractive gravity.
\end{enumerate}
The NEC is not derived from self-modeling; it is an input to the sign
chain.  The tensor equation~\eqref{eq:einstein} itself holds
independently of the NEC.

\subsubsection{Parameter identification}
\label{sec:parameters}

Newton's constant is set by the UV entropy density:
$\Gnewton = 1/(4\eta)$.  From the lattice-to-continuum mapping
(Sec.~\ref{sec:continuum}), $\eta = \eta_{\mathrm{lattice}}/a^{d-1}$,
where $\eta_{\mathrm{lattice}} \leq \log n$ is the entropy per
boundary bond (from the channel capacity
bound~\eqref{eq:channel-bound}).  This gives
\begin{equation}
  \Gnewton \geq \frac{a^{d-1}}{4\log n},
  \label{eq:g-bound}
\end{equation}
identifying the lattice spacing $a$ with the Planck length up to a
factor of $\sqrt{\log n}$.  In $d = 3$ spatial dimensions:
$a \lesssim 2\ell_P\sqrt{\log n}$, where
$\ell_P = \sqrt{\Gnewton}$.

The Lieb--Robinson velocity $\vlr$ (Eq.~\ref{eq:vlr}) maps to the
speed of light~$c$ in the continuum limit.  The cosmological
constant~$\Lambda$ is undetermined; this paper does not predict its
value.

\subsection{Route B: Lovelock uniqueness ($d \geq 2$)}
\label{sec:route-b}

Route~A derives Einstein's equation by explicit computation of
$\delta\ent_{\mathrm{UV}}$ and $\delta\ent_{\mathrm{mat}}$ using the
Raychaudhuri equation and the conformal modular Hamiltonian.  The
conformal approximation (Sec.~\ref{sec:conformal-approx}) limits this
route to the conformal regime, which is controlled in $d = 1$ but
problematic in $d \geq 2$ where the infrared theory likely has N\'eel
order rather than conformal symmetry.

Route~B circumvents this limitation entirely.  It uses only three
ingredients, none of which require a conformal modular Hamiltonian:

\begin{enumerate}
\item \emph{Emergent geometry from entanglement.}  The self-modeling
  lattice defines an emergent geometry via the thermal time
  identification (Sec.~\ref{sec:equilibrium}): the modular flow of the
  lattice state defines the geometric flow, and the entanglement
  structure encodes the metric.  This holds in all spatial dimensions.
\item \emph{Entanglement first law.}  Perturbations of the entanglement
  entropy satisfy $\delta\ent = \delta\langle\modH\rangle$
  (Eq.~\ref{eq:first-law}), an exact identity.  Combined with the
  area-law structure (Sec.~\ref{sec:arealaw}), this implies that
  perturbations of entanglement are equivalent to perturbations of the
  emergent geometry.  We assume the resulting equation relating metric
  perturbations to matter perturbations is a symmetric 2-tensor equation
  built from the metric and at most second derivatives (\emph{tensoriality
  assumption}).  Symmetry follows from the symmetry of the metric;
  the second-derivative restriction follows from the entanglement first
  law being first-order in the perturbation, which couples to the
  modular Hamiltonian (a local operator at leading order in the
  derivative expansion).  Divergence-freeness is then a consequence
  of Lovelock's theorem rather than an input: Lovelock identifies the
  unique tensor, and that tensor happens to be divergence-free
  (the contracted Bianchi identity).
\item \emph{Lovelock's theorem.}  Lovelock~\cite{Lovelock1971,
  Lovelock1972} proved that in $d + 1$ spacetime dimensions with
  $d \geq 2$, the \emph{unique} symmetric, divergence-free 2-tensor
  constructed from the metric $\gab$ and its first and second
  derivatives is a linear combination of $\Gab$ and $\gab$:
  \begin{equation}
    \mathcal{E}_{ab} = \alpha\, \Gab + \beta\, \gab
  \end{equation}
  for constants $\alpha$, $\beta$.  (In $d + 1 > 4$, Lovelock gravity
  allows higher-curvature terms built from the generalized
  Euler densities; in $d + 1 = 3$ and $d + 1 = 4$, the Einstein
  tensor and cosmological term are the unique possibilities.)
\end{enumerate}

Combining these: since the entanglement-geometry relation is a
symmetric 2-tensor equation with at most second derivatives in the
emergent $(d{+}1)$-dimensional spacetime, Lovelock's theorem constrains
the geometric side to be
\begin{equation}
  \alpha\, \Gab + \beta\, \gab = \text{(matter source)},
  \label{eq:einstein-lovelock}
\end{equation}
in $d + 1 \leq 4$ dimensions.  (In $d + 1 > 4$, higher Lovelock terms
such as the Gauss--Bonnet combination are permitted; restricting to
$d + 1 \leq 4$, or equivalently to at most second-order field equations,
excludes them.)  Identifying $\alpha = 1$, $\beta = \Lambda$, and
$\Gnewton = 1/(4\eta)$ by matching the UV entanglement density
(Sec.~\ref{sec:parameters}) recovers Eq.~\eqref{eq:einstein}.

Route~B constrains the \emph{geometric sector}: it shows the left side
must be Einsteinian.  It does not independently derive the matter
coupling $8\pi\Gnewton\,\Tab$ on the right side; that identification
comes from Route~A's explicit modular Hamiltonian computation, or more
generally from identifying the modular Hamiltonian perturbation with
the local stress tensor (the entanglement first law).  The two routes
are therefore genuinely complementary: Route~A supplies the matter
coupling, Route~B supplies the geometric uniqueness.

\emph{Complementarity of the two routes.}---The Jacobson route
(Route~A) derives the \emph{specific form} of the Einstein equation by
explicit computation, but requires the conformal modular Hamiltonian
and is therefore controlled only in the conformal regime.  The Lovelock
route (Route~B) derives the \emph{uniqueness} of the Einstein equation
from general principles, but requires the assumption that the
entanglement-geometry relation produces a tensorial equation with at
most second-order derivatives.

\begin{table}
\caption{Comparison of the two routes to Einstein's equations.
\label{tab:routes}}
\begin{ruledtabular}
\begin{tabular}{lll}
  & Route A (Jacobson) & Route B (Lovelock) \\ \hline
Strength & Explicit derivation & Uniqueness argument \\
Best in & $d = 1$ (WZW CFT) & $d \geq 2$ (all) \\
Requires & Conformal $\modH$ & Tensorial equation \\
Weakness & Conformal approx.\ & No explicit $\Tab$ \\
\end{tabular}
\end{ruledtabular}
\end{table}

In $d = 1$ ($1{+}1$ spacetime), $\Gab = 0$ identically for all metrics,
so both routes produce a trivially satisfied equation. Route~A gives
$\Lambda\,\gab = 8\pi\Gnewton\,\Tab$ (a cosmological-constant equation)
via the exact WZW modular Hamiltonian; Lovelock's theorem provides no
constraint.  The $d = 1$ case serves as a controlled consistency check
rather than a demonstration of nontrivial gravitational dynamics.

In $d \geq 2$, Route~B forces the geometric sector to be Einsteinian
by uniqueness, without requiring conformal symmetry of the infrared
theory.  Route~A provides the explicit matter coupling in the conformal
window $a \ll \ell \ll 1/m$.  Neither route is assumption-free: Route~A
needs the conformal approximation, Route~B needs the tensoriality
assumption.  Their complementarity strengthens the overall case, with
each route addressing the other's principal weakness.
